\documentclass[aspectratio=169]{beamer}
\usetheme{Madrid}
\usecolortheme{default}

\usepackage{amsmath,amssymb,amsfonts,bm,mathtools}
\usepackage{physics}
\usepackage{braket}

\title{Lipkin Hamiltonian to Pauli Strings}
\subtitle{Quasispin Representation, Binary Encoding, and Direct Qubit Mapping}
\author{Morten Hjorth-Jensen}
\date{Spring 2026}

\begin{document}

\frame{\titlepage}

\begin{frame}{Goal}
We want to rewrite the Lipkin Hamiltonian
\[
H = \frac{\epsilon}{2}\left(J_z-\frac{N}{2}\right)
+ \frac{V}{2}\left(J_+^2+J_-^2\right)
+ \frac{W}{2}J_z^2
\]
in terms of tensor products of Pauli operators \(I,X,Y,Z\).

\medskip

Two cases will be treated:
\begin{enumerate}
\item \(N=4\) particles in the quasispin \(j=2\) sector, encoded on \textbf{3 qubits} (compressed binary encoding).
\item \(N=6\) particles represented on \textbf{6 qubits} (direct collective-spin realization).
\end{enumerate}

The emphasis is on \emph{how} the mapping is constructed, not only on the final formulas.
\end{frame}

\begin{frame}{The quasispin algebra}
The Lipkin model is naturally written in terms of collective quasispin operators
\[
J_x,\qquad J_y,\qquad J_z,
\]
with ladder operators
\[
J_\pm = J_x \pm iJ_y.
\]

They satisfy the \(su(2)\) commutation relations
\[
[J_x,J_y]=iJ_z,\qquad [J_y,J_z]=iJ_x,\qquad [J_z,J_x]=iJ_y,
\]
or equivalently
\[
[J_z,J_\pm]=\pm J_\pm,\qquad [J_+,J_-]=2J_z.
\]
\end{frame}

\begin{frame}{Useful algebraic simplification}
A key identity is
\[
J_+^2 + J_-^2
=
(J_x+iJ_y)^2 + (J_x-iJ_y)^2
=
2(J_x^2-J_y^2).
\]

Therefore the Hamiltonian may also be written as
\[
H
=
\frac{\epsilon}{2}\left(J_z-\frac{N}{2}\right)
+
V\left(J_x^2-J_y^2\right)
+
\frac{W}{2}J_z^2.
\]

This form is particularly convenient for direct qubit mappings.
\end{frame}

\begin{frame}{Quasispin basis}
In the fixed-\(j\) representation, the natural basis is
\[
\ket{j,m},\qquad m=-j,-j+1,\dots,j.
\]

The standard actions are
\[
J_z\ket{j,m}=m\ket{j,m},
\]
\[
J_+\ket{j,m}
=
\sqrt{j(j+1)-m(m+1)}\,\ket{j,m+1},
\]
\[
J_-\ket{j,m}
=
\sqrt{j(j+1)-m(m-1)}\,\ket{j,m-1}.
\]

Thus \(J_z\) is diagonal, while \(J_\pm\) shift \(m\) by one unit.
\end{frame}

\begin{frame}{Strategy for producing Pauli strings}
There are two standard routes:
\begin{enumerate}
\item \textbf{Compressed encoding:} encode the quasispin basis \(\ket{j,m}\) directly into computational basis states of a smaller qubit register.
\item \textbf{Direct qubit mapping:} express the collective operators as sums over physical qubit Pauli operators and expand the Hamiltonian.
\end{enumerate}

In these notes:
\begin{itemize}
\item the \(N=4\) case uses route 1,
\item the \(N=6\) case uses route 2.
\end{itemize}
\end{frame}

\section{Part I: Four particles encoded on three qubits}

\begin{frame}{Case \(N=4\): quasispin dimension}
For \(N=4\), the quasispin is
\[
j = \frac{N}{2} = 2.
\]

Hence the Hilbert space dimension in the \(j=2\) sector is
\[
2j+1 = 5.
\]

Since
\[
2^2=4 < 5 \le 8 = 2^3,
\]
the smallest binary encoding requires \textbf{3 qubits}.
\end{frame}

\begin{frame}{Choice of binary encoding for \(N=4\)}
We choose the map
\[
\ket{j=2,m=-2} \leftrightarrow \ket{000},
\]
\[
\ket{j=2,m=-1} \leftrightarrow \ket{001},
\]
\[
\ket{j=2,m=0} \leftrightarrow \ket{010},
\]
\[
\ket{j=2,m=1} \leftrightarrow \ket{011},
\]
\[
\ket{j=2,m=2} \leftrightarrow \ket{100}.
\]

The remaining computational states
\[
\ket{101},\ket{110},\ket{111}
\]
are unphysical and should be excluded or penalized.
\end{frame}

\begin{frame}{Projectors in terms of Pauli \(Z\)}
For one qubit,
\[
\ket{0}\bra{0} = \frac{I+Z}{2},
\qquad
\ket{1}\bra{1} = \frac{I-Z}{2}.
\]

Thus for a three-qubit computational basis state \(\ket{abc}\),
\[
\ket{abc}\bra{abc}
=
\frac{1}{8}(I+(-1)^a Z_2)(I+(-1)^b Z_1)(I+(-1)^c Z_0).
\]

This gives a systematic way to build diagonal operators such as \(J_z\) and \(J_z^2\).
\end{frame}

\begin{frame}{Transition operators in terms of \(X\) and \(Y\)}
For one qubit,
\[
\ket{1}\bra{0} = \frac{X-iY}{2},
\qquad
\ket{0}\bra{1} = \frac{X+iY}{2}.
\]

Therefore an off-diagonal transition such as
\[
\ket{a_2a_1a_0}\bra{b_2b_1b_0}
\]
is written as a tensor product of:
\begin{itemize}
\item projectors \(\frac{I\pm Z}{2}\) on unchanged bits,
\item ladder pieces \(\frac{X\mp iY}{2}\) on flipped bits.
\end{itemize}

This is the basic rule for expanding \(J_+\), \(J_-\), and \(J_+^2+J_-^2\).
\end{frame}

\begin{frame}{Matrix form of \(J_z\) for \(j=2\)}
In the encoded basis, \(J_z\) is diagonal:
\[
J_z
=
\sum_{m=-2}^{2} m \ket{m}\bra{m}.
\]

With the chosen binary map, this becomes
\[
J_z
=
-2\ket{000}\bra{000}
-\ket{001}\bra{001}
+0\ket{010}\bra{010}
+\ket{011}\bra{011}
+2\ket{100}\bra{100}.
\]

We take the operator to vanish on the three unused basis states.
\end{frame}

\begin{frame}{Explicit Pauli decomposition of \(J_z\) (3 qubits)}
Expanding the projectors gives
\[
J_z
=
-\frac{1}{4} IZI
+\frac{1}{4} IZZ
-\frac{1}{2} ZII
-\frac{1}{2} ZIZ
-\frac{3}{4} ZZI
-\frac{1}{4} ZZZ.
\]

Here the notation means, for example,
\[
IZZ = I\otimes Z\otimes Z,
\qquad
ZII = Z\otimes I\otimes I.
\]

This formula is exact on the chosen 5-dimensional encoded subspace.
\end{frame}

\begin{frame}{Explicit Pauli decomposition of \(J_z^2\) (3 qubits)}
Since
\[
J_z^2 = \sum_{m=-2}^{2} m^2 \ket{m}\bra{m},
\]
we obtain
\[
J_z^2
=
\frac{5}{4} III
+\frac{3}{4} IIZ
+IZI
+IZZ
+\frac{1}{4} ZII
-\frac{1}{4} ZIZ.
\]

This again acts exactly on the chosen encoded subspace and is set to zero on unused basis states.
\end{frame}

\begin{frame}{Action of the ladder operators for \(j=2\)}
For \(j=2\),
\[
J_+\ket{-2}=2\ket{-1},
\qquad
J_+\ket{-1}=\sqrt{6}\ket{0},
\]
\[
J_+\ket{0}=\sqrt{6}\ket{1},
\qquad
J_+\ket{1}=2\ket{2},
\qquad
J_+\ket{2}=0.
\]

Thus
\[
J_+
=
2\ket{001}\bra{000}
+\sqrt{6}\ket{010}\bra{001}
+\sqrt{6}\ket{011}\bra{010}
+2\ket{100}\bra{011}.
\]

And \(J_- = J_+^\dagger\).
\end{frame}

\begin{frame}{Why \(J_+^2+J_-^2\) is simpler than \(J_+\) itself}
Because \(J_+^2\) changes \(m\) by two units, it connects only
\[
\ket{-2}\leftrightarrow \ket{0},
\qquad
\ket{-1}\leftrightarrow \ket{1},
\qquad
\ket{0}\leftrightarrow \ket{2}.
\]

For \(j=2\),
\[
J_+^2\ket{-2}=2\sqrt{6}\ket{0},
\]
\[
J_+^2\ket{-1}=6\ket{1},
\]
\[
J_+^2\ket{0}=2\sqrt{6}\ket{2}.
\]

Therefore
\[
J_+^2+J_-^2
\]
is real and Hermitian, and its Pauli decomposition contains only real coefficients.
\end{frame}

\begin{frame}{Explicit \(J_+^2+J_-^2\) decomposition for \(N=4\)}
The exact three-qubit decomposition is
\[
J_+^2+J_-^2
=
\frac{3+\sqrt{6}}{2}\,IXI
+
\frac{\sqrt{6}-3}{2}\,IXZ
+
\frac{\sqrt{6}}{2}\,XXI
+
\frac{\sqrt{6}}{2}\,XXZ
\]
\[
\qquad
+
\frac{\sqrt{6}}{2}\,YYI
+
\frac{\sqrt{6}}{2}\,YYZ
+
\frac{3+\sqrt{6}}{2}\,ZXI
+
\frac{\sqrt{6}-3}{2}\,ZXZ.
\]

This is the key off-diagonal part of the Lipkin Hamiltonian in the compressed 3-qubit encoding.
\end{frame}

\begin{frame}{Final \(N=4\), 3-qubit Hamiltonian}
Using
\[
H = \frac{\epsilon}{2}(J_z-2) + \frac{V}{2}(J_+^2+J_-^2) + \frac{W}{2}J_z^2,
\]
the encoded Hamiltonian is
\[
H_{N=4}^{(3q)}
=
\left(-\epsilon+\frac{5W}{8}\right)III
+\frac{3W}{8}IIZ
+\left(-\frac{\epsilon}{8}+\frac{W}{2}\right)IZI
+\left(\frac{\epsilon}{8}+\frac{W}{2}\right)IZZ
\]
\[
+\left(-\frac{\epsilon}{4}+\frac{W}{8}\right)ZII
+\left(-\frac{\epsilon}{4}-\frac{W}{8}\right)ZIZ
-\frac{3\epsilon}{8}ZZI
-\frac{\epsilon}{8}ZZZ
\]
\[
+\frac{V(3+\sqrt{6})}{4}IXI
+\frac{V(\sqrt{6}-3)}{4}IXZ
+\frac{V\sqrt{6}}{4}XXI
+\frac{V\sqrt{6}}{4}XXZ
\]
\[
+\frac{V\sqrt{6}}{4}YYI
+\frac{V\sqrt{6}}{4}YYZ
+\frac{V(3+\sqrt{6})}{4}ZXI
+\frac{V(\sqrt{6}-3)}{4}ZXZ.
\]
\end{frame}

\begin{frame}{Physical-subspace penalty for the 3-qubit encoding}
Because only five of the eight computational basis states are physical, it is often useful to add a penalty
\[
H_{\mathrm{pen}} = \lambda P_{\mathrm{unphys}},
\]
where
\[
P_{\mathrm{unphys}}
=
\ket{101}\bra{101}
+\ket{110}\bra{110}
+\ket{111}\bra{111}.
\]

This suppresses leakage into unused basis states.
Each projector can again be written in terms of \(I\) and \(Z\) only.
\end{frame}

\begin{frame}{Interpretation of the 3-qubit mapping}
The compressed mapping is efficient because:
\begin{itemize}
\item the \(j=2\) quasispin sector has dimension only \(5\),
\item binary encoding uses only \(\lceil \log_2 5\rceil = 3\) qubits,
\item all quasispin operators become finite sums of Pauli strings.
\end{itemize}

The price is that one must carefully enforce the physical code space.
\end{frame}

\section{Part II: Six particles on six qubits}

\begin{frame}{Case \(N=6\): direct collective-spin realization}
For \(N=6\), one natural choice is to use one qubit per particle.
Then the collective quasispin operators are represented directly as
\[
J_x = \frac{1}{2}\sum_{i=1}^{6} X_i,
\qquad
J_y = \frac{1}{2}\sum_{i=1}^{6} Y_i,
\qquad
J_z = \frac{1}{2}\sum_{i=1}^{6} Z_i,
\]
where \(X_i,Y_i,Z_i\) act on qubit \(i\).
\end{frame}

\begin{frame}{Collective spin squares}
Because \(X_i^2=Y_i^2=Z_i^2=I\),
\[
J_x^2
=
\frac{1}{4}\sum_{i=1}^{6} I
+
\frac{1}{4}\sum_{i\neq j} X_iX_j
=
\frac{6}{4}I + \frac{1}{2}\sum_{i<j} X_iX_j.
\]

Similarly,
\[
J_y^2
=
\frac{6}{4}I + \frac{1}{2}\sum_{i<j} Y_iY_j,
\]
\[
J_z^2
=
\frac{6}{4}I + \frac{1}{2}\sum_{i<j} Z_iZ_j.
\]
\end{frame}

\begin{frame}{Rewrite the interaction term}
Using
\[
J_+^2+J_-^2 = 2(J_x^2-J_y^2),
\]
we get
\[
\frac{V}{2}(J_+^2+J_-^2)
=
V(J_x^2-J_y^2).
\]

Substituting the previous expansions,
\[
J_x^2-J_y^2
=
\frac{1}{2}\sum_{i<j}\left(X_iX_j-Y_iY_j\right),
\]
since the identity contributions cancel.
Therefore
\[
\frac{V}{2}(J_+^2+J_-^2)
=
\frac{V}{2}\sum_{i<j}\left(X_iX_j-Y_iY_j\right).
\]
\end{frame}

\begin{frame}{Rewrite the linear \(J_z\) term}
We have
\[
\frac{\epsilon}{2}\left(J_z-\frac{N}{2}\right)
=
\frac{\epsilon}{2}\left(\frac{1}{2}\sum_{i=1}^{6} Z_i - 3I\right).
\]

Hence
\[
\frac{\epsilon}{2}\left(J_z-\frac{N}{2}\right)
=
\frac{\epsilon}{4}\sum_{i=1}^{6} Z_i
-
\frac{3\epsilon}{2}I.
\]
\end{frame}

\begin{frame}{Rewrite the \(J_z^2\) term}
Using
\[
J_z^2
=
\frac{6}{4}I + \frac{1}{2}\sum_{i<j} Z_iZ_j,
\]
we obtain
\[
\frac{W}{2}J_z^2
=
\frac{3W}{4}I
+
\frac{W}{4}\sum_{i<j} Z_iZ_j.
\]
\end{frame}

\begin{frame}{Final \(N=6\), 6-qubit Hamiltonian}
Collecting all pieces,
\[
H_{N=6}^{(6q)}
=
\left(-\frac{3\epsilon}{2}+\frac{3W}{4}\right)I
+
\frac{\epsilon}{4}\sum_{i=1}^{6} Z_i
+
\frac{V}{2}\sum_{i<j}\left(X_iX_j-Y_iY_j\right)
+
\frac{W}{4}\sum_{i<j} Z_iZ_j.
\]

This is the clean direct Pauli-string representation on 6 qubits.
\end{frame}

\begin{frame}{How many Pauli terms appear for \(N=6\)?}
For 6 qubits:
\begin{itemize}
\item one identity term,
\item \(6\) one-body \(Z_i\) terms,
\item \(\binom{6}{2}=15\) two-body \(X_iX_j\) terms,
\item \(15\) two-body \(Y_iY_j\) terms,
\item \(15\) two-body \(Z_iZ_j\) terms.
\end{itemize}

So the total number of Pauli strings is
\[
1 + 6 + 15 + 15 + 15 = 52.
\]
This is still highly structured, since all coefficients come from only three parameters \(\epsilon,V,W\).
\end{frame}

\begin{frame}{Explicit pairwise form for the interaction term}
The interaction contribution is
\[
\frac{V}{2}
\Big(
X_1X_2-X_1Y_2
\Big)
\]
No---that is not the correct structure.

The correct two-body interaction is
\[
\frac{V}{2}\sum_{1\le i<j\le 6}
\left(X_iX_j-Y_iY_j\right).
\]

Expanded explicitly, it includes pairs such as
\[
X_1X_2-Y_1Y_2,\quad
X_1X_3-Y_1Y_3,\quad \dots,\quad
X_5X_6-Y_5Y_6.
\]

This pairwise all-to-all form is a hallmark of the direct collective-spin mapping.
\end{frame}

\begin{frame}{Direct mapping versus compressed mapping}
For \(N=6\), the direct 6-qubit representation has advantages:
\begin{itemize}
\item extremely simple operator formulas,
\item immediate interpretation in terms of collective spin,
\item no leakage into unphysical binary code states.
\end{itemize}

Its disadvantage is qubit cost:
\begin{itemize}
\item the physical symmetric subspace has dimension only \(N+1=7\),
\item but the full register has dimension \(2^6=64\).
\end{itemize}

So this is simple but not qubit-optimal.
\end{frame}

\section{General Lessons and Comparison}

\begin{frame}{Comparison of the two mappings}
\begin{center}
\begin{tabular}{lcc}
\hline
Case & Qubits used & Main advantage \\
\hline
\(N=4\), compressed & 3 & minimal qubit count \\
\(N=6\), direct & 6 & simplest Pauli expansion \\
\hline
\end{tabular}
\end{center}

More conceptually:
\begin{itemize}
\item compressed encodings reduce qubit count but require code-space management,
\item direct collective encodings use more qubits but produce transparent many-body Pauli Hamiltonians.
\end{itemize}
\end{frame}

\begin{frame}{General recipe for quasispin-to-Pauli conversion}
To convert a Lipkin quasispin Hamiltonian to Pauli strings:

\begin{enumerate}
\item Choose an encoding of the \(\ket{j,m}\) basis into computational basis states.
\item Express diagonal operators (\(J_z\), \(J_z^2\)) using computational basis projectors.
\item Express off-diagonal operators (\(J_\pm\), \(J_\pm^2\)) using transition operators built from \(X\) and \(Y\).
\item Simplify using
\[
J_+^2+J_-^2 = 2(J_x^2-J_y^2)
\]
whenever a direct collective-spin realization is available.
\item Collect coefficients into sums of Pauli strings.
\end{enumerate}
\end{frame}

\begin{frame}{Summary}
We have shown how to transform the Lipkin Hamiltonian
\[
H = \frac{\epsilon}{2}\left(J_z-\frac{N}{2}\right)
+ \frac{V}{2}(J_+^2+J_-^2)
+ \frac{W}{2}J_z^2
\]
into Pauli-string form in two concrete cases:

\begin{itemize}
\item \(N=4\): exact \textbf{3-qubit compressed binary encoding} of the \(j=2\) quasispin sector,
\item \(N=6\): exact \textbf{6-qubit direct collective-spin mapping}.
\end{itemize}

The central ideas are:
\begin{itemize}
\item diagonal operators from \(Z\)-projectors,
\item transitions from \(X\) and \(Y\),
\item collective-spin identities for simplification.
\end{itemize}
\end{frame}

\begin{frame}{End}
\centering
{\Large Questions?}
\end{frame}

\end{document}
